\cleardoublepage
\thispagestyle{empty}
\bigtitle{IV die infra oct. Sanct. Cord. Jesu.}
\addcontentsline{toc}{section}{XV Junii} \phantomsection
\semiduplex

\header{xv junii 2026}

\office{ad primam}
\addcontentsline{toc}{subsection}{Ad Primam} \phantomsection

\pateravecredo

\gscore[]{℣.}{Deus_in_adiutorium_Tonus_ferialis_no_laus_tibi}

\gscore[Hy.]{3.}{hy_jam_lucis_orto_sidere_in_festo_s._cordis_jesu_solesmes_1961}

\gscore[Ant.]{1 f.}{unus_militum}
\psalmus{53}{53_1f}{53_1}
\psalmus{118.\textsc{i}}{118_1_1f}{118_1_1}
\psalmus{118.\textsc{ii}}{118_2_1f}{118_2_1}
\smallscore{an_unus_militum}

\capitulum{1 Tim. 1, 17}
\lettrine{R}{e}gi sæculórum immortáli et invisíbili,~† soli Deo honor et glória~* in sǽcula
sæculórum.

\rr Deo grátias.

\gscore[]{Resp. brev.}{rb_christe_fili_dei_vivi_sc}

\smallscore{versiculus_ad_primam}

  \vv Dóminus vobíscum. \rr Et cum spíritu tuo.
  
  \oratio
  
  Orémus.
  
 \lettrine{D}{o}mine Deus omnípotens, qui ad princípium huius diéi nos perveníre
fecísti: tua nos hódie salva virtúte; ut in hac die ad nullum declinémus peccátum,
sed semper ad tuam justítiam faciéndam nostra procédant elóquia, dirigántur
cogitatiónes et ópera. Per Dóminum.

\rr Amen.

\vv Dóminus vobíscum. \rr Et cum spíritu tuo.

\smallscore{BD_ad_completorium}

\smalltitle{Martyrologium.}

\begin{enpars}
\lettrine[lraise=0.2]{J}{u}ne 16\textsuperscript{th} 2026, the first day of the Moon, were born into the better life:
At Besançon, in Gaul, the holy martyrs Ferreolus a Priest and Ferrutio a Deacon, who were sent forth by the blessed Irenaeus, Bishop (of Lyon,) to preach the Word of God, whereafter they were diverse ways tormented and then beheaded under Claudius the judge.
At Tarsus, in Cilicia, under the Emperor Diocletian, the holy martyrs Quiricus and his mother Julitta.
Quiricus was a little boy of three years old when his mother was horribly scourged before the President Alexander, he vehemently wept and lamented her, whereupon he was killed by being dashed against the steps of the judgment-seat. Julitta, after horrid stripes and grievous torments, was beheaded, and so finished her testimony.
At Mainz, the holy martyrs Aureus and his sister Justina, and the others who were at Communion in the Church when they were massacred by the Huns, who were wasting Germany, (in the year 451).
At Limasol, in Cyprus, holy Tycho, Bishop (of that see,) in the time of the Emperor Theodosius the younger.
At Lyon, (in the year 551,) the blessed Aurelian, Bishop of Arles.
At Nantes, in Brittany, (in the year 310,) the holy Confessor Similian, Bishop (of that see.)
At Meissen, in Germany, (in the year 1106,) holy Benno, Bishop (of that see.)
At the village of La Louvesc, in the Diocese of Vienne, in the Ardèche, (in the year 1640,) the holy Confessor John Francis Regis, of the Society of Jesus, a man of wonderful love and long-suffering in seeking the salvation of souls, whom Pope Clement \scspace{xii} added to the list of the Saints.
In Brabant, (in the year 1246,) the holy Virgin Lutgard.
\end{enpars}

\vv Et álibi aliórum plurimórum san\-ctórum Mártyrum et Confessórum, atque san\-ctárum Vírginum.

\rr Deo grátias.

\vv Pretiósa in conspéctu Dómini. \rr Mors Sanctórum ejus.

Sancta María et omnes Sancti intercédant pro nobis ad Dóminum, ut nos
mereámur ab eo adjuvári et salvári, qui vivit et regnat in sǽcula sæculórum.

\rr Amen.

\vv Deus in adjutórium meum inténde.

\rr Dómine, ad adjuvándum me festína.

\vv Deus in adjutórium meum inténde.

\rr Dómine, ad adjuvándum me festína.

\vv Deus in adjutórium meum inténde.

\rr Dómine, ad adjuvándum me festína.

\vv Glória Patri, et Fílio,~* et Spirítui Sancto.

\rr Sicut erat in princípio, et nunc, et semper,~* et in sǽcula sæculórum. Amen.

\vv Kýrie eléison.

\rr Christe eléison. Kýrie eléison.

\vv Pater noster \textit{secreto.}

\vv Et ne nos indúcas in tentatiónem:

\rr Sed líbera nos a malo.

\vv Réspice in servos tuos, Dómine, et in ópera tua, et dírige fílios eórum.

\rr Et sit splendor Dómini Dei nostri super nos, et ópera mánuum nostrárum
dírige super nos, et opus mánuum nostrárum dírige.

\vv Glória Patri, et Fílio,~* et Spirítui Sancto.

\rr Sicut erat in princípio, et nunc, et semper,~* et in sǽcula sæculórum. Amen.

Orémus.

\lettrine{D}{i}rígere et sanctificáre, régere et gubernáre dignáre, Dómine Deus, Rex
cæli et terræ, hódie corda et córpora nostra, sensus, sermónes et a\-ctus nostros in
lege tua, et in opéribus mandatórum tuórum: ut hic et in ætérnum, te auxiliánte,
salvi et líberi esse mereámur, Salvátor mundi: Qui vivis et regnas in sǽcula
sæculórum.

\rr Amen.

\vv Jube, domne, benedícere.

\smalltitle{Benedictio.}

Dies et actus nostros in sua pace dispónat Dóminus omnípotens.

\rr Amen.

\lectiobrevis{Eph. 3, 17 -- 19}
\lettrine{I}{n} caritáte radicáti et fundáti, ut possítis comprehéndere cum ómnibus sanctis, quæ sit latitúdo et longitúdo, et sublímitas et profúndum; scire étiam supereminéntem sciéntiæ caritátem Christi, ut impleámini in omnem plenitúdinem Dei.

\vv Tu autem, Dómine, miserére nobis.

\rr Deo grátias.

\vv Adjutórium nostrum in nómine Dómini.

\rr Qui fecit cælum et terram.

\vv Benedícite.

\rr Deus.

\lettrine{D}{o}minus nos benedícat, \cc{} et ab omni malo deféndat, et ad vitam perdúcat ætérnam. Et fidélium ánimæ per misericórdiam Dei requiéscant in pace.

\rr Amen.

\secreto

\office{ad missam}
\addcontentsline{toc}{subsection}{Ad Missam} \phantomsection
\proper{Introitus}{Ps. 32, 11, 19; 1}

\texteparacol{\lettrine{C}{o}gitatiónes Cordis ejus in generatióne et generatiónem: ut éruat a morte ánimas eórum et alat eos in fame.
℣. Exsultáte, justi, in Dómino: rectos decet collaudátio.
Glória Patri. Sicut erat.
Cogitatiónes.}{english}{\noindent The thoughts of his heart are to all generations: to deliver them from death and preserve them in spite of famine. ℣. Exult, you just, in the Lord; praise from the upright is fitting. Glory be to the Father. As it was in the beginning. The thoughts.}

\masspart{Oratio}

\texteparacol{\lettrine{D}{e}us, qui nobis in Corde Fílii tui, nostris vulneráto peccátis, infinítos dile\-ctiónis thesáuros misericórditer largíri dignáris:~† concéde quǽsumus; ut illi devótum pietátis nostræ præstántes obséquium,~* dignæ quoque satisfa\-ctiónis exhibeámus offícium. Per eúmdem Dóminum.

\lettrine[ante=2]{D}{a} Ecclésiæ tuæ, quǽsumus Dómine, san\-ctis Martýribus tuis Vito, Modésto atque Crescéntia intercedéntibus, supérbe non sápere, sed tibi plácita humilitáte profícere:~† ut prava despíciens,~* quæcúmque re\-cta sunt, líbera exérceat caritáte.

\lettrine[ante=3]{C}{o}ncéde nos famúlos tuos, quæsumus Domíne Deus, perpétua mentis et corpóris sanitáte gaudére:~† et gloriósa beátæ Mariæ semper Virgínis intercessióne,~* a præsénti liberári tristítia, et æterna perfrúi lætítia. Per Dóminum.}{english}{\noindent O God, who hast suffered the heart of thy Son to be wounded by our sins, and in that very heart hast bestowed on us the abundant riches of thy love: grant that the devout homage of our hearts, which we render unto him, may of thy mercy be deemed a recompence acceptable in thy sight. Through the same Lord.

\noindent 2. Grant o Lord, we pray, that through the intercession of thy holy martyrs Vitus, Modestus and Crescentia, thy church be not high-minded, but grow in that humility which is pleasing to thee, so that despising evil, she may practice whatever is good with unbounded charity.

\noindent 3. Grant us, thy servants, o Lord God, we beseech thee to enjoy lasting health of mind and body; and by the intercession of glorious and blessed Mary, ever Virgin, may we be delivered from present sorrow and partake to the full of eternal happiness. Through our Lord.}

\reading{Epistola}{Eph. 3, 8 -- 12, 14 -- 19}

\texteparacol{\lettrine{F}{r}atres: Mihi, ómnium san\-ctórum mínimo, data est grátia hæc, in géntibus evangelizáre investigábiles divítias Christi, et illumináre o\-mnes, quæ sit dispensátio sacraménti abscónditi a sǽculis in Deo, qui ó\-mnia creávit: ut innotéscat principátibus et potestátibus in cæléstibus per Ecclésiam multifórmis sapiéntia Dei, secúndum præfinitiónem sæculórum, quam fecit in Christo Jesu, Dómino nostro, in quo habémus fidúciam et accéssum in confidéntia per fidem ejus. Hujus rei grátia fle\-cto génua mea ad Patrem Dómini nostri Jesu Christi, ex quo omnis patérnitas in cælis ei in terra nominátur, ut det vobis, secúndum divítias glóriæ suæ, virtúte corroborári per Spíritum ejus in interiórem hóminem, Christum habitáre per fidem in córdibus vestris: in caritáte radicáti et fundáti, ut póssitis comprehéndere cum ó\-mnibus san\--ctis, quæ sit latitúdo, et longitúdo, et sublímitas, et profúndum: scire étiam supereminéntem sciéntiæ caritátem Christi, ut impleámini in omnem plenitúdinem Dei.}{english}{\noindent Brethren: To me, the very least of all saints, there was given this grace, to announce among the Gentiles the good tidings of the unfathomable riches of Christ, and to enlighten all men as to what is the dispensation of the mystery which has been hidden from eternity in God, Who created all things; in order that through the Church there be made known to the Principalities and the Powers in the heavens the manifold wisdom of God according to the eternal purpose which He accomplished in Christ Jesus our Lord. In Him we have assurance and confident access through faith in Him. For this reason I bend my knees to the Father of our Lord Jesus Christ, from Whom all fatherhood in heaven and on earth receives its name, that He may grant you from His glorious riches to be strengthened with power through His Spirit unto the progress of the inner man; and to have Christ dwelling through faith in your hearts: so that, being rooted and grounded in love, you may be able to comprehend with all the saints what is the breadth and length and height and depth, and to know Christ’s love which surpasses knowledge, in order that you may be filled unto all the fullness of God.}

\proper{Graduale \& Alleluia}{Ps. 24, 8 -- 9; Matt. 11, 29}

\texteparacol{\lettrine{D}{u}lcis et rectus Dóminus: propter hoc legem dabit delinquéntibus in via.
℣. Díriget mansúetos in judício, docébit mites vias suas.

\lettrine{A}{l}elúia, allelúia. Tóllite jugum meum super vos, et díscite a me, quia mitis sum et húmilis Corde, et inveniétis réquiem animábus vestris. Allelúia.}{english}{\noindent Good and upright is the Lord; thus He shows sinners the way.
℣. He guides the humble to justice; he teaches the humble his way.

\noindent Alleluia, alleluia. Take my yoke upon you, and learn from me, for I am meek, and humble of heart: and you will find rest for your souls. Alleluia.}

\reading{Evangelium}{Joann. 19, 31 -- 37}

\texteparacol{\lettrine{I}{n} illo témpore: Judǽi, quóniam Parascéve erat, ut non remanérent in cruce córpora sábbato, erat enim magnus dies ille sábbati, rogavérunt Pilátum, ut frangeréntur eórum crura, et tolleréntur. Venérunt ergo mílites: et primi quidem fregérunt crura et alteríus, qui crucifíxus est cum eo. Ad Jesum autem cum veníssent, ut vidérunt eum jam mórtuum, non fregérunt ejus crura, sed unus mílitum láncea latus ejus apéruit, et contínuo exívit sanguis et aqua. Et qui vidit, testimónium perhíbuit: et verum est testimónium ejus. Et ille scit quia vera dicit, ut et vos credátis. Fa\-cta sunt enim hæc ut Scriptúra implerétur: Os non comminuétis ex eo. Et íterum alia Scri\-ptúra dicit: Vidébunt in quem transfixérunt.}{english}{\noindent At that time, the Jews, since it was the Preparation Day, in order that the bodies might not remain upon the cross on the Sabbath,  for that Sabbath was a solemn day, besought Pilate that their legs might be broken, and that they might be taken away. The soldiers therefore came and broke the legs of the first, and of the other, who had been crucified with him. But when they came to Jesus, and saw that he was already dead, they did not break his legs; but one of the soldiers opened his side with a lance, and immediately there came out blood and water. And he who saw it has borne witness, and his witness is true; and he knows that he tells the truth, that you also may believe. For these things came to pass that the Scripture might be fulfilled: ``Not a bone of him shall you break.'' And again another Scripture says: ``They shall look upon him whom they have pierced.''}

\proper{Offertorium}{Ps. 68, 21}

\texteparacol{\lettrine{I}{m}propérium exs\-pectávi Cor meum et misériam: et sustínui, qui simul mecum contristarétur, et non fuit: consolántem me quæsívi, et non invéni.}{english}{\noindent My heart expected reproach and misery; I looked for sympathy, but there was none; and for comforters, and I found none.}

\masspart{Secreta}
\texteparacol{\lettrine{R}{e}spice, quǽsumus, Dómine, ad ineffábilem Cordis dilé\-cti Fílii tui caritátem: ut quod offérimus sit tibi munus accé\-ptum et nostrórum expiátio deli\-ctórum.
Per eúmdem Dóminum.

\lettrine[ante=2]{S}{i}cut glóriam divínæ poténtiæ múnera pro San\-ctis obláta testántur: sic nobis effé\-ctum, Dómine, tuæ salvatiónis impéndant.

\lettrine[ante=3]{T}{u}a, Dómine, propitiatióne, et beátæ Maríæ semper Vírginis intercessióne, ad perpétuam atque præséntem hæc oblátio nobis profíciat prosperitátem et pacem.
Per Dóminum.}{english}{\noindent Look, we beseech thee, o Lord, upon the heart of thy beloved Son, with its boundless love, so that what we offer, may be an acceptable gift in thy sight and an atonement for our sins.
Through the same Lord.

\noindent 2. As the gifts, we offer in honor of thy saints, o Lord, bear witness to the glory of thy divine power; so may they win for us the benefits of thy salvation.

\noindent 3. Through thy mercy, o Lord, and the intercession of blessed Mary ever Virgin, may this oblation obtain for us prosperity and peace both now and forever.
Through our Lord.}

\masspart{Præfatio de sacratissimo Cordis Jesu}

\texteparacol{\lettrine{V}{e}re dignum et justum est, æquum et salutáre, nos tibi semper et ubíque grátias ágere: Dómine san\-cte, Pater omnípotens, ætérne Deus: Qui Unigénitum tuum, in Cruce pendéntem, láncea mílitis transfígi voluísti: ut apértum Cor, divínæ largitátis sacrárium, torréntes nobis fúnderet miseratiónis et grátiæ: et, quod amóre nostri flagráre numquam déstitit, piis esset réquies et pæniténtibus pater et salútis refúgium. Et ídeo cum Angelis et Archángelis, cum Thronis et Dominatiónibus cumque omni milítia cæléstis exércitus hymnum glóriæ tuæ cánimus, sine fine dicéntes:}{english}{\noindent It is truly meet and just, right and for our salvation, that we should at all times, and in all places, give thanks unto Thee, O holy Lord, Father almighty, everlasting God; Who didst will that Thine only-begotten Son, while hanging on the cross, should be pierced by a soldier's spear, that the Heart thus opened, a shrine of divine bounty, should pour out on us streams of mercy and grace, and that what never ceased to burn with love for us, should be a resting-place to the devout, and open as a refuge of salvation to the penitent. And therefore with angels and archangels, with thrones and dominions, and with all the hosts of the heavenly army, we sing the hymn of thy glory, evermore saying:}

\proper{Communio}{Joann. 19, 34}

\texteparacol{\lettrine{U}{n}us mílitum láncea latus ejus apéruit, et contínuo exívit sanguis et aqua.}{english}{\noindent One of the soldiers opened his side with a lance, and immediately there came out blood and water.}

\masspart{Postcommunio}
\texteparacol{\lettrine{P}{r}ǽbeant nobis, Dómine Jesu, divínum tua san\-cta fervórem: quo dulcíssimi Cordis tui suavitáte percé\-pta; discámus terréna despícere, et amáre cæléstia:
Qui vivis et regnas.

\lettrine[ante=2]{R}{e}pléti, Dómine, benedi\-ctióne sollémni: quǽsumus; ut, per intercessiónem san\-ctórum Mártyrum tuórum Viti, Modésti et Crescéntiæ, medicína sacraménti et corpóribus nostris prosit et méntibus.

\lettrine[ante=3]{S}{u}mptis, Dómine, salútis nostræ subsídiis: da, quǽsumus, beátæ Maríæ semper Vírginis patrocíniis nos ubíque prótegi; in cujus veneratióne hæc tuæ obtúlimus majestáti. Per Dóminum.}{english}{\noindent May thy sacrament, o Lord Jesus, give us holy zeal, so that, seeing the sweetness of thy most loving heart, we may learn to despise the things of earth and love those of heaven.
Who livest and reignest.

\noindent 2. We, who have been filled with thy solemn blessing, ask, o Lord, that, through the intercession of thy holy martyrs Vitus, Modestus and Crescentia, this healing sacrament may benefit both our bodies and our souls.

\noindent 3. Grant, o Lord, that we, who had partaken of the aids of salvation, may be everywhere protected by the patronage of blessed Mary ever Virgin, in whose honor we have offered these gifts of thy majesty. Through our Lord.}

\office{ad vesperas}
\addcontentsline{toc}{subsection}{Ad  Vesperas} \phantomsection
%\label{in_2_vesperis_sc}\phantomsection
\label{vespers_2_sd}\phantomsection

\gscore[]{℣.}{Deus_In_Adjutorium_Festivus}

\gscore[1. Ant.]{1 f.}{unus_militum}
\psalmus{109}{109_1f}{109_1}
\smallscore{an_unus_militum}

\gscore[2. Ant.]{7 c.}{stans_jesus} %% LaTeX's tracking is HORRIBLE with this psalm. Lots of white for no reason. %%needs fixing
\psalmus{110}{110_7c}{110_7}
\smallscore{an_stans_jesus}

\gscore[3. Ant.]{3. a2}{in_caritate}
\psalmus{115}{115_3a2}{115_3a2_g}
\smallscore{an_in_caritate}
%\smallskip

\gscore[4. Ant.]{4. E}{venite}
\psalmus{127}{127_4E}{127_4E}
\smallscore{an_venite_ad_me}

\gscore[5. Ant.]{5. a}{fili}
\psalmus{147}{147_5}{147_5}
\smallscore{an_fili_praebe}
%an_ignem

\capitulum{Eph. 3, 8 – 9}

\lettrine{F}{r}ratres: Mihi ó\-mnium san\-ctórum mínimo data est grátia hæc,~† in géntibus evangelizáre investigábiles divítias Christi:~* et illumináre omnes, quæ sit dispensátio sacraménti abscónditi a sǽculis in Deo.

\rr Deo grátias.

\hymnus
   \gscore[]{3.}{hy_en_ut_superba_solesmes_1961}
   
   \smallscore{versiculus_haurietis_tonus_solius}

   \admagnificat
   
   \gscore[]{1.f}{ad_jesum_autem}
    
    \smallscore{Magnificat_1f}
    
  \magnificattexte{Magnificat_1_ordinaire}
  
  \smallscore{an_ad_jesum_autem}
  
  \vv Dóminus vobíscum. \rr Et cum spíritu tuo.
  
  \oratio
  
  Orémus.
  
  \lettrine{D}{e}us, qui nobis in Corde Fílii tui, nostris vulneráto peccátis, infinítos dile\-ctiónis thesáuros misericórditer largíri dignáris:~† concéde quǽsumus; ut illi devótum pietátis nostræ præstántes obséquium,~* dignæ quoque satisfa\-ctiónis exhibeámus offícium. Per eúmdem Dóminum.
  
  \rr Amen.

   \vv Dóminus vobíscum. \rr Et cum spíritu tuo.
   
   \gscore[]{2.}{BD_Semiduplex_Vesp}
   
  \vv Fidélium ánimæ per misericórdiam Dei requiéscant in pace.
  
  \rr Amen.
  
 \rubrique{Tunc dicitur Completorium. Si non dicendum:}
  
  \secreto
  
\vv Dóminus det nobis suam pacem.
  
\rr Et vitam ætérnam. Amen.

\rubrique{\normaltext{Salve Regína,} etc., \normaltext{\pageref{an_salve_regina_solesmes}.}}